\section{RF moment transfer and dissipative balance}
For one travelling mode with $n=|k|$, chirality $\sigma=\mathrm{sgn}(k)$, frequency ratio $\nu_j$ and $P=\alpha^2n^2|a_jc_k|^2r^{2n-2}$, the zero-orbit-speed linear response gives
\begin{equation}
\overline M_{\rm RF,total}=\sigma nP\frac{\lambda\nu_j}{D},\qquad
\overline M_{\rm slow}=\sigma(n-1)P\frac{\lambda\nu_j}{D},\quad
D=(\nu_j^2+\alpha\beta)^2+(\lambda\nu_j)^2.
\end{equation}
The difference is the RF moment assigned to micromotion; drag on micromotion must be subtracted to obtain the slow-orbit transfer. For $k=+2$ and $-4$, slow torques scale as $+A^2r^2$ and $-B^2\nu r^6$. Their competition plus radius-dependent localization permits different signs for different $q/m$. Drag balances the injected torque; the system is driven-dissipative and no minimization of total angular momentum is implied. In the tested pair, counterrotation persists when Coulomb is set to zero, so Coulomb is not the source of the signs.
